\section{Выводы}
Реализованы наивный поиск подстроки и алгоритм Кнута — Морриса — Пратта со счётчиками
сравнений, их корректность подтверждена сверкой с \texttt{str.find} на 300 случайных входах и
граничными случаями.

Замеры подтвердили теоретические оценки. На худшем входе время наивного поиска растёт как
$n^2$, а КМП — линейно, и при $n = 16~000$ разница достигает 779,3 против 2,115~мс. На
случайном тексте оба алгоритма линейны и почти равны по времени, поэтому преимущество КМП
проявляется только на текстах с длинными частичными совпадениями образца.
